% ============================================================
%  GUÍA 6: Trigonometría Analítica y Gráficas
%  Curso de Introducción a la Universidad — Precálculo y Cálculo I
%  Autor: Ing. Gabriel Astudillo
%  Clase cubierta: 6
% ============================================================

\documentclass[12pt, a4paper]{article}

% ── Paquetes ─────────────────────────────────────────────────
\usepackage{mathwizards-guia}
\mwguia{Guía 6 — Trigonometría Analítica y Gráficas}{Ing. Gabriel Astudillo $\cdot$ Curso Preuniversitario}

\begin{document}

% ── Portada / Título ─────────────────────────────────────────
\mwportada{Guía de Estudio 6}{Trigonometría Analítica y Gráficas}{Grados, Radianes, Longitud de Arco, Circunferencia Unitaria, Ondas e Identidades}

\vspace{4pt}
\begin{cuadroinfo}
\small
\textbf{Presentación:} Esta guía cubre la trigonometría necesaria para el cálculo: medición angular en grados y radianes, longitud de arco, la circunferencia unitaria, los ángulos notables y de referencia, las gráficas de las funciones trigonométricas —amplitud, período y desfase— y las identidades fundamentales. Con teoría, ejemplos resueltos y propuestos con clave.
\\[4pt]
\textbf{Nivel de Dificultad:}\quad
\facil{} Básico / Directo \quad
\medio{} Intermedio \quad
\dificil{} Desafiante \quad
\muyDificil{} Muy Difícil / Reto
\end{cuadroinfo}

\vspace{8pt}

  \section{Sistema de Medición Angular y Circunferencia Unitaria}
  % ============================================================

  \subsection{Grados y radianes}

  \begin{cuadroteoria}
    \textbf{Conversión:}
    $$180^\circ = \pi \text{ rad} \qquad\Longrightarrow\qquad \text{radianes} = \frac{\pi}{180} \times \text{grados}$$
    $$\text{grados} = \frac{180}{\pi} \times \text{radianes}$$

    \textbf{Longitud de arco:} $s = r\theta$ (con $\theta$ en radianes).

    \textbf{Área de un sector circular:} $A = \frac{1}{2}r^2 \theta$ (con $\theta$ en radianes).
  \end{cuadroteoria}

  \begin{cuadrosolucion}
    \textbf{Ángulos notables en la circunferencia unitaria} ($r = 1$):

    \begin{center}
      \begin{tabular}{lcccccc}
        \toprule
        \textbf{Ángulo $\theta$} & $0$       & $\dfrac{\pi}{6}$      & $\dfrac{\pi}{4}$      & $\dfrac{\pi}{3}$      & $\dfrac{\pi}{2}$ & $\pi$       \\\\
        \textbf{Grados}          & $0^\circ$ & $30^\circ$            & $45^\circ$            & $60^\circ$            & $90^\circ$       & $180^\circ$ \\\\
        \midrule
        $\sin\theta$             & $0$       & $\dfrac{1}{2}$        & $\dfrac{\sqrt{2}}{2}$ & $\dfrac{\sqrt{3}}{2}$ & $1$              & $0$         \\\\
        $\cos\theta$             & $1$       & $\dfrac{\sqrt{3}}{2}$ & $\dfrac{\sqrt{2}}{2}$ & $\dfrac{1}{2}$        & $0$              & $-1$        \\\\
        \bottomrule
      \end{tabular}
    \end{center}
  \end{cuadrosolucion}

  \begin{cuadroadvertencia}
    \textbf{Coordenadas en la circunferencia unitaria:} para un ángulo $\theta$, el punto sobre la circunferencia unitaria es:
    $$(\cos\theta,\ \sin\theta).$$
    El seno es la coordenada $y$ y el coseno es la coordenada $x$.
  \end{cuadroadvertencia}

  \subsection{Ejercicios resueltos}

  \textbf{Ejemplo R1.} Convierte $120^\circ$ a radianes.

  \begin{cuadrosolucion}
    \textbf{Solución:} $120^\circ \times \dfrac{\pi}{180} = \dfrac{120\pi}{180} = \dfrac{2\pi}{3}$ rad.
  \end{cuadrosolucion}

  \textbf{Ejemplo R2.} Convierte $\dfrac{5\pi}{4}$ a grados.

  \begin{cuadrosolucion}
    \textbf{Solución:} $\dfrac{5\pi}{4} \times \dfrac{180}{\pi} = \dfrac{5 \cdot 180}{4} = 225^\circ$.
  \end{cuadrosolucion}

  \textbf{Ejemplo R3.} Un círculo tiene radio $6$ cm. Encuentra la longitud de arco para un ángulo central de $\dfrac{\pi}{3}$ radianes y el área del sector correspondiente.

  \begin{cuadrosolucion}
    \textbf{Solución:}
    \begin{itemize}
      \item Longitud: $s = r\theta = 6 \cdot \dfrac{\pi}{3} = 2\pi \approx 6{,}28$ cm.
      \item Área: $A = \dfrac{1}{2}r^2\theta = \dfrac{1}{2}(36)\left(\dfrac{\pi}{3}\right) = 6\pi \approx 18{,}85$ cm$^2$.
    \end{itemize}
  \end{cuadrosolucion}

  \textbf{Ejemplo R4.} Determina $\sin\left(\frac{2\pi}{3}\right)$ y $\cos\left(\frac{2\pi}{3}\right)$ usando la circunferencia unitaria.

  \begin{cuadrosolucion}
    \textbf{Solución:} $\frac{2\pi}{3} = 120^\circ$, en el segundo cuadrante. El ángulo de referencia es $60^\circ$.
    \begin{itemize}
      \item $\cos\left(\frac{2\pi}{3}\right) = -\cos\left(\frac{\pi}{3}\right) = -\dfrac{1}{2}$ (negativo en el 2do cuadrante).
      \item $\sin\left(\frac{2\pi}{3}\right) = \sin\left(\frac{\pi}{3}\right) = \dfrac{\sqrt{3}}{2}$ (positivo).
    \end{itemize}
  \end{cuadrosolucion}

  \textbf{Ejemplo R5.} Una rueda de radio $0{,}5$ m gira un ángulo de $3$ radianes. ¿Qué distancia recorre un punto en el borde?

  \begin{cuadrosolucion}
    \textbf{Solución:} $s = r\theta = 0{,}5 \cdot 3 = 1{,}5$ m.
  \end{cuadrosolucion}

  \newpage

  \subsection{Ejercicios propuestos}

  \begin{enumerate}[leftmargin=*]
    \item \facil{} Convierte a radianes: $30^\circ$, $45^\circ$, $90^\circ$, $180^\circ$, $270^\circ$.
    \item \facil{} Convierte a grados: $\dfrac{\pi}{6}$, $\dfrac{\pi}{4}$, $\dfrac{\pi}{3}$, $\dfrac{3\pi}{2}$.
    \item \facil{} Calcula: $\sin\left(\frac{\pi}{6}\right)$, $\cos\left(\frac{\pi}{4}\right)$, $\tan\left(\frac{\pi}{4}\right)$.
    \item \facil{} Calcula: $\sin(0)$, $\cos(0)$, $\sin\left(\frac{\pi}{2}\right)$, $\cos(\pi)$.
    \item \medio{} Un círculo de radio $10$ cm tiene un ángulo central de $72^\circ$. Convierte el ángulo a radianes y calcula la longitud del arco.
    \item \medio{} Calcula el área de un sector circular de radio $5$ m y ángulo $\dfrac{\pi}{4}$ rad.
    \item \medio{} Determina $\sin\left(\frac{5\pi}{6}\right)$, $\cos\left(\frac{5\pi}{6}\right)$ usando el ángulo de referencia.
    \item \medio{} Determina $\sin\left(\frac{7\pi}{4}\right)$, $\cos\left(\frac{7\pi}{4}\right)$.
    \item \dificil{} Encuentra todos los ángulos $\theta$ en $[0, 2\pi)$ tales que $\sin\theta = \dfrac{1}{2}$.
    \item \dificil{} Encuentra todos los ángulos $\theta$ en $[0, 2\pi)$ tales que $\cos\theta = -\dfrac{\sqrt{2}}{2}$.
    \item \dificil{} Una bicicleta tiene ruedas de radio $35$ cm. Si la rueda gira $200$ vueltas, ¿qué distancia recorre la bicicleta en metros?
    \item \dificil{} Un péndulo de $2$ m de largo oscila con un ángulo de $10^\circ$. ¿Qué longitud de arco recorre la punta del péndulo en una oscilación completa (ida y vuelta)?
  \end{enumerate}

  \newpage

  % ============================================================

  \section{Gráficas de Funciones Trigonométricas e Identidades Básicas}
  % ============================================================

  \subsection{Parámetros de una onda senoidal}

  \begin{cuadroteoria}
    \textbf{Forma general:}
    $$y = A\sin(Bx - C) + D \qquad \text{o} \qquad y = A\cos(Bx - C) + D$$

    \begin{itemize}[leftmargin=*]
      \item \textbf{Amplitud:} $|A|$ (la mitad de la distancia entre el máximo y el mínimo).
      \item \textbf{Período:} $\dfrac{2\pi}{|B|}$ (espacio horizontal de un ciclo completo).
      \item \textbf{Desfase:} $\dfrac{C}{B}$ (desplazamiento horizontal; a la derecha si $C > 0$).
      \item \textbf{Desplazamiento vertical:} $D$ (hacia arriba si $D > 0$).
    \end{itemize}

    \textbf{Función tangente:} $y = \tan x = \dfrac{\sin x}{\cos x}$. Período $\pi$; asíntotas verticales donde $\cos x = 0$.
  \end{cuadroteoria}

  \begin{cuadrosolucion}
    \textbf{Identidades pitagóricas:}
    \begin{align*}
      \sin^2 x + \cos^2 x & = 1        \\
      1 + \tan^2 x        & = \sec^2 x \\
      1 + \cot^2 x        & = \csc^2 x
    \end{align*}

    \textbf{Otras identidades básicas:}
    $$\tan x = \frac{\sin x}{\cos x}, \qquad \cot x = \frac{\cos x}{\sin x}, \qquad \sec x = \frac{1}{\cos x}, \qquad \csc x = \frac{1}{\sin x}$$

    \textbf{Identidad del ángulo doble (anticipo):}
    $$\sin(2\theta) = 2\sin\theta\cos\theta$$
    Esta identidad se usará más adelante al derivar funciones trigonométricas (Guía 6).
  \end{cuadrosolucion}

  \begin{cuadroinfo}
    \textbf{Nota:} Las funciones trigonométricas inversas ($\arcsin$, $\arccos$, $\arctan$) se estudiarán en la Guía 6, donde aprenderás a derivarlas.
  \end{cuadroinfo}

  \subsection{Ejercicios resueltos}

  \textbf{Ejemplo R1.} Determina amplitud, período, desfase y desplazamiento vertical de $y = 2\sin(3x - \pi) + 1$.

  \begin{cuadrosolucion}
    \textbf{Solución:} Comparando con $A\sin(Bx - C) + D$:
    \begin{itemize}
      \item Amplitud: $|2| = 2$.
      \item Período: $\dfrac{2\pi}{3}$.
      \item Desfase: $\dfrac{C}{B} = \dfrac{\pi}{3}$ (hacia la derecha).
      \item Desplazamiento vertical: $+1$.
    \end{itemize}
  \end{cuadrosolucion}

  \textbf{Ejemplo R2.} Grafica $y = \cos x$ y $y = \cos x + 2$ en el mismo sistema.

  \begin{cuadrosolucion}
    \begin{center}
      \begin{tikzpicture}[scale=0.9]
        \draw[thin, gray!40] (-3.5,-1.5) grid (7.5,4.5);
        \draw[thick, -Latex] (-4,0) -- (8,0) node[right] {$x$};
        \draw[thick, -Latex] (0,-2) -- (0,5) node[above] {$y$};
        \draw[domain=-3.5:7, smooth, very thick, gray!70, dashed] plot (\x, {cos(deg(\x))});
        \draw[domain=-3.5:7, smooth, very thick, mwprimario] plot (\x, {cos(deg(\x))+2});
        \draw[gray] (-1.5,3) node[above] {$y = \cos x$};
        \draw[gray] (4,3.5) node[above] {$y = \cos x + 2$};
        \draw (-pi,0) node[below left] {$\pi$};
        \draw (pi,0) node[below left] {$\pi$};
      \end{tikzpicture}
    \end{center}
  \end{cuadrosolucion}

  \textbf{Ejemplo R3.} Verifica la identidad $\sin x \cdot \tan x = \dfrac{1 - \cos^2 x}{\cos x}$.

  \begin{cuadrosolucion}
    \textbf{Solución:} Partimos del lado izquierdo:
    $$\sin x \cdot \tan x = \sin x \cdot \frac{\sin x}{\cos x} = \frac{\sin^2 x}{\cos x} = \frac{1 - \cos^2 x}{\cos x},$$
    usando $\sin^2 x + \cos^2 x = 1$. Se verifica.
  \end{cuadrosolucion}

  \textbf{Ejemplo R4.} Escribe la ecuación de una onda senoidal con amplitud 3, período $\pi$, desfase $\dfrac{\pi}{4}$ a la derecha y desplazamiento vertical $-2$.

\begin{cuadrosolucion}
  \textbf{Solución:} Período $= \dfrac{2\pi}{B} = \pi \Rightarrow B = 2$. Desfase $= \dfrac{C}{B} = \dfrac{\pi}{4} \Rightarrow \dfrac{C}{2} = \dfrac{\pi}{4} \Rightarrow C = \dfrac{\pi}{2}$.
  $$y = 3\sin\left(2x - \frac{\pi}{2}\right) - 2.$$
\end{cuadrosolucion}

\newpage

\subsection{Ejercicios propuestos}

\begin{enumerate}[leftmargin=*, resume]
  \item \facil{} Determina amplitud y período de $y = 4\sin(2x)$.
  \item \facil{} Determina amplitud y período de $y = \dfrac{1}{2}\cos(3x)$.
  \item \facil{} Determina desfase y desplazamiento vertical de $y = \sin\left(x - \dfrac{\pi}{2}\right) + 3$.
  \item \facil{} Calcula $\tan\left(\frac{\pi}{4}\right)$ usando $\sin$ y $\cos$.
  \item \medio{} Determina todos los parámetros (amplitud, período, desfase, desplazamiento vertical) de $y = -2\cos\left(\dfrac{x}{2} + \dfrac{\pi}{3}\right) + 1$. (Cuidado: factoriza $B$ antes de hallar el desfase.)
  \item \medio{} Simplifica: $\dfrac{\sin^2 x}{1 - \cos x}$ (pista: usa $\sin^2 x = 1 - \cos^2 x$ y factoriza).
  \item \medio{} Verifica la identidad: $\cos x \tan x = \sin x$.
  \item \medio{} Verifica la identidad: $\dfrac{1}{\sin x} - \sin x = \dfrac{\cos^2 x}{\sin x}$.
  \item \dificil{} Escribe la ecuación de una onda coseno con amplitud 2, período $4\pi$, desfase $\dfrac{\pi}{3}$ a la izquierda y desplazamiento vertical $5$.
  \item \dificil{} Encuentra todos los valores de $x$ en $[0, 2\pi)$ que satisfacen $\sin^2 x = \dfrac{1}{4}$.
  \item \dificil{} Demuestra la identidad: $\dfrac{1 + \tan^2 x}{\sec x} = \sec x$.
  \item \dificil{} La altura del agua en un puerto está modelada por $h(t) = 3 + 2\sin\left(\dfrac{\pi t}{6}\right)$, donde $t$ está en horas y $h$ en metros.
        \begin{enumerate}[label=\alph*)]
          \item ¿Cuál es la altura máxima y mínima del agua?
          \item ¿Cuál es el período de la marea?
          \item ¿En qué instantes la altura es $3$ m en las primeras 12 horas?
        \end{enumerate}
\end{enumerate}

\newpage

% ============================================================

% ──────────────────────────────────────────────────────────
%  NUEVOS EJERCICIOS PROPUESTOS (26) — INSERCIÓN MANUAL
% ──────────────────────────────────────────────────────────
\subsection{Ejercicios propuestos: Ángulos, Arcos y Circunferencia Unitaria}

\begin{enumerate}[leftmargin=*, resume]
  \item \facil{} Convierte $120^\circ$ a radianes.

  \item \facil{} Convierte $\dfrac{5\pi}{6}$ a grados.

  \item \medio{} Calcula la longitud del arco subtendido por un ángulo central de $\dfrac{\pi}{4}$ en una circunferencia de radio 8 cm.

  \item \medio{} Evalúa exactamente: $\sin \dfrac{\pi}{6}$, $\cos \dfrac{\pi}{3}$ y $\tan \dfrac{\pi}{4}$.

  \item \medio{} Determina el ángulo de referencia de $225^\circ$ y usa la circunferencia unitaria para hallar $\sin 225^\circ$ y $\cos 225^\circ$.

  \item \dificil{} Evalúa exactamente $\sin \dfrac{5\pi}{3}$ y $\cos \dfrac{5\pi}{3}$.

  \item \dificil{} Un ángulo central de $2{,}5$ rad subtiende un arco de 15 cm. Determina el radio de la circunferencia.

  \item \dificil{} Si $\cos \theta = \dfrac{3}{5}$ y $\theta$ está en el cuarto cuadrante, halla $\sin \theta$ y $\tan \theta$.

  \item \muyDificil{} Una rueda de 30 cm de radio gira un ángulo de $240^\circ$. Determina la distancia recorrida por un punto del borde.
\end{enumerate}

\subsection{Ejercicios propuestos: Gráficas e Identidades}

\begin{enumerate}[leftmargin=*, resume]
  \item \facil{} Determina el período de $y = \sin(2x)$.

  \item \facil{} Determina la amplitud y el período de $y = 3\cos\left(\dfrac{x}{2}\right)$.

  \item \medio{} Determina el desfase de $y = \sin\left(x - \dfrac{\pi}{3}\right)$ respecto a $y = \sin x$.

  \item \medio{} Escribe la ecuación de una función senoidal con amplitud 2, período $\pi$ y desfase 0.

  \item \medio{} Simplifica las expresiones: $\dfrac{\sin x}{\cos x}$ \quad y \quad $\sin^2 x + \cos^2 x$.

  \item \dificil{} Simplifica usando la identidad pitagórica:
        $$\frac{1 - \cos^2 x}{\sin x}$$

  \item \dificil{} Simplifica la expresión trigonométrica:
        $$\tan x \cdot \cos x$$

  \item \dificil{} Determina el valor máximo y el valor mínimo de $y = 2 + 3\sin x$.

  \item \muyDificil{} Demuestra la identidad:
        $$\frac{\sin x}{1 + \cos x} = \frac{1 - \cos x}{\sin x}$$
\end{enumerate}

\subsection{Ejercicios propuestos: Retos Integradores}

\begin{enumerate}[leftmargin=*, resume]
  \item \facil{} Convierte $-135^\circ$ a radianes.

  \item \medio{} Evalúa exactamente: $\sec \dfrac{\pi}{3}$ \quad y \quad $\csc \dfrac{\pi}{6}$.

  \item \medio{} Halla todos los ángulos $\theta \in [0, 2\pi)$ tales que $\sin \theta = \dfrac{1}{2}$.

  \item \medio{} Determina el período de $y = \tan(3x)$.

  \item \dificil{} \textbf{Modelado:} La altura del agua en un puerto sigue $h(t) = 5 + 3\cos\left(\dfrac{\pi t}{6}\right)$ con $t$ en horas. Determina la altura máxima, la mínima y el período de la marea.

  \item \dificil{} Simplifica la suma de fracciones trigonométricas:
        $$\frac{\cos x}{1 - \sin x} + \frac{\cos x}{1 + \sin x}$$

  \item \muyDificil{} Demuestra la identidad:
        $$\frac{1}{\tan x + \cot x} = \sin x \cos x$$

  \item \muyDificil{} Un péndulo de 1{,}2 m oscila con una amplitud angular de $\dfrac{\pi}{12}$ rad. Determina la longitud del arco que barre su extremo.
\end{enumerate}

% ============================================================
\section{Clave de Respuestas}
% ============================================================

\subsection*{Sección 1: Ángulos y Circunferencia Unitaria (Ejercicios 1--12)}

\begin{enumerate}[leftmargin=*, label=\textbf{\arabic*.}]
  \item $\dfrac{\pi}{6}$, $\dfrac{\pi}{4}$, $\dfrac{\pi}{2}$, $\pi$, $\dfrac{3\pi}{2}$
  \item $30^\circ$, $45^\circ$, $60^\circ$, $270^\circ$
  \item $\sin\frac{\pi}{6} = \frac{1}{2}$; $\cos\frac{\pi}{4} = \frac{\sqrt{2}}{2}$; $\tan\frac{\pi}{4} = 1$
  \item $0$, $1$, $1$, $-1$
  \item $72^\circ = \dfrac{2\pi}{5}$ rad. $s = 10 \cdot \dfrac{2\pi}{5} = 4\pi \approx 12{,}57$ cm.
  \item $A = \dfrac{1}{2}(25)\left(\dfrac{\pi}{4}\right) = \dfrac{25\pi}{8} \approx 9{,}82$ m$^2$.
  \item Ángulo de referencia $\frac{\pi}{6}$. $\sin\frac{5\pi}{6} = \frac{1}{2}$; $\cos\frac{5\pi}{6} = -\frac{\sqrt{3}}{2}$.
  \item Ángulo de referencia $\frac{\pi}{4}$. $\sin\frac{7\pi}{4} = -\frac{\sqrt{2}}{2}$; $\cos\frac{7\pi}{4} = \frac{\sqrt{2}}{2}$.
  \item $\theta = \dfrac{\pi}{6}$ y $\theta = \dfrac{5\pi}{6}$.
  \item $\theta = \dfrac{3\pi}{4}$ y $\theta = \dfrac{5\pi}{4}$.
  \item $s = 200 \cdot 2\pi r = 200 \cdot 2\pi (0{,}35) = 140\pi \approx 439{,}8$ m.
  \item $10^\circ = \dfrac{\pi}{18}$ rad. Ida: $s = 2 \cdot \dfrac{\pi}{18} = \dfrac{\pi}{9} \approx 0{,}349$ m. Oscilación completa (ida y vuelta) $= 2 \cdot \dfrac{\pi}{9} = \dfrac{2\pi}{9} \approx 0{,}698$ m.
\end{enumerate}


\subsection*{Sección 2: Gráficas Trigonométricas e Identidades (Ejercicios 13--24)}

\begin{enumerate}[leftmargin=*, label=\textbf{\arabic*.}]
  \setcounter{enumi}{12}
  \item Amplitud $4$; período $\dfrac{2\pi}{2} = \pi$.
  \item Amplitud $\dfrac{1}{2}$; período $\dfrac{2\pi}{3}$.
  \item Desfase $\dfrac{\pi}{2}$ a la derecha; desplazamiento vertical $+3$.
  \item $\tan\frac{\pi}{4} = \dfrac{\sin\frac{\pi}{4}}{\cos\frac{\pi}{4}} = \dfrac{\frac{\sqrt{2}}{2}}{\frac{\sqrt{2}}{2}} = 1$.
  \item Reescribimos: $y = -2\cos\left[\dfrac{1}{2}\left(x + \dfrac{2\pi}{3}\right)\right] + 1$. Amplitud $2$; período $\dfrac{2\pi}{1/2} = 4\pi$; desfase $-\dfrac{2\pi}{3}$ (a la izquierda); desplazamiento vertical $+1$.
  \item $\dfrac{\sin^2 x}{1 - \cos x} = \dfrac{1 - \cos^2 x}{1 - \cos x} = \dfrac{(1 - \cos x)(1 + \cos x)}{1 - \cos x} = 1 + \cos x$.
  \item $\cos x \tan x = \cos x \cdot \dfrac{\sin x}{\cos x} = \sin x$. Se verifica.
  \item $\dfrac{1}{\sin x} - \sin x = \dfrac{1 - \sin^2 x}{\sin x} = \dfrac{\cos^2 x}{\sin x}$. Se verifica.
  \item Período $4\pi \Rightarrow B = \dfrac{2\pi}{4\pi} = \dfrac{1}{2}$. Desfase a la izquierda $\dfrac{\pi}{3}$: $\dfrac{C}{B} = -\dfrac{\pi}{3} \Rightarrow C = -\dfrac{\pi}{6}$. Ecuación: $y = 2\cos\left(\dfrac{x}{2} + \dfrac{\pi}{6}\right) + 5$.
  \item $\sin^2 x = \frac{1}{4} \Rightarrow \sin x = \pm \frac{1}{2}$. En $[0, 2\pi)$: $x = \dfrac{\pi}{6}, \dfrac{5\pi}{6}, \dfrac{7\pi}{6}, \dfrac{11\pi}{6}$.
  \item $1 + \tan^2 x = \sec^2 x$. Entonces $\dfrac{1 + \tan^2 x}{\sec x} = \dfrac{\sec^2 x}{\sec x} = \sec x$. Se verifica.
  \item a) Altura máxima: $3 + 2 = 5$ m; mínima: $3 - 2 = 1$ m. \quad b) Período: $\dfrac{2\pi}{\pi/6} = 12$ horas. \quad c) $h = 3$ cuando $\sin\left(\frac{\pi t}{6}\right) = 0$, es decir $\frac{\pi t}{6} = 0, \pi, 2\pi$ $\Rightarrow$ $t = 0, 6, 12$ (en las primeras 12 horas).
\end{enumerate}


% ──────────────────────────────────────────────────────────
%  NUEVAS RESPUESTAS (26) — INSERCIÓN MANUAL
% ──────────────────────────────────────────────────────────
\subsection*{Sección 4: Ángulos, Arcos y Circunferencia Unitaria (Ejercicios 25--33)}
\begin{enumerate}[leftmargin=*, label=\textbf{\arabic*.}]
  \setcounter{enumi}{24}
  \item $\dfrac{2\pi}{3}$ rad
  \item $150^\circ$
  \item $s = 8 \cdot \dfrac{\pi}{4} = 2\pi$ cm $\approx 6{,}28$ cm
  \item $\dfrac{1}{2}$, $\dfrac{1}{2}$, $1$
  \item Ángulo de referencia $45^\circ$; tercer cuadrante: $\sin 225^\circ = -\dfrac{\sqrt{2}}{2}$, $\cos 225^\circ = -\dfrac{\sqrt{2}}{2}$
  \item $\sin \dfrac{5\pi}{3} = -\dfrac{\sqrt{3}}{2}$; \quad $\cos \dfrac{5\pi}{3} = \dfrac{1}{2}$
  \item $r = \dfrac{s}{\theta} = \dfrac{15}{2{,}5} = 6$ cm
  \item $\sin \theta = -\dfrac{4}{5}$; \quad $\tan \theta = -\dfrac{4}{3}$
  \item $240^\circ = \dfrac{4\pi}{3}$ rad; \quad $s = 30 \cdot \dfrac{4\pi}{3} = 40\pi$ cm $\approx 125{,}66$ cm
\end{enumerate}

\subsection*{Sección 5: Gráficas e Identidades (Ejercicios 34--42)}
\begin{enumerate}[leftmargin=*, label=\textbf{\arabic*.}]
  \setcounter{enumi}{33}
  \item Período $\pi$
  \item Amplitud 3; período $4\pi$
  \item Desfase $\dfrac{\pi}{3}$ hacia la derecha
  \item $y = 2\sin(2x)$
  \item $\dfrac{\sin x}{\cos x} = \tan x$; \quad $\sin^2 x + \cos^2 x = 1$
  \item $\dfrac{\sin^2 x}{\sin x} = \sin x$
  \item $\tan x \cos x = \dfrac{\sin x}{\cos x} \cos x = \sin x$
  \item Máximo $2 + 3 = 5$; \quad mínimo $2 - 3 = -1$
  \item Multiplicando por el conjugado: $\dfrac{\sin x}{1 + \cos x} \cdot \dfrac{1 - \cos x}{1 - \cos x} = \dfrac{\sin x(1 - \cos x)}{1 - \cos^2 x} = \dfrac{1 - \cos x}{\sin x}$.
\end{enumerate}

\subsection*{Sección 6: Retos Integradores (Ejercicios 43--50)}
\begin{enumerate}[leftmargin=*, label=\textbf{\arabic*.}]
  \setcounter{enumi}{42}
  \item $-\dfrac{3\pi}{4}$ rad
  \item $\sec \dfrac{\pi}{3} = 2$; \quad $\csc \dfrac{\pi}{6} = 2$
  \item $\theta = \dfrac{\pi}{6}$ o $\theta = \dfrac{5\pi}{6}$
  \item Período $\dfrac{\pi}{3}$
  \item Máximo $5 + 3 = 8$ m; \quad mínimo $5 - 3 = 2$ m; \quad período $\dfrac{2\pi}{\pi/6} = 12$ horas
  \item $\dfrac{\cos x(1 + \sin x) + \cos x(1 - \sin x)}{1 - \sin^2 x} = \dfrac{2\cos x}{\cos^2 x} = 2\sec x$
  \item $\dfrac{1}{\frac{\sin x}{\cos x} + \frac{\cos x}{\sin x}} = \dfrac{\sin x \cos x}{\sin^2 x + \cos^2 x} = \sin x \cos x$.
  \item $s = 1{,}2 \cdot \dfrac{\pi}{12} = 0{,}1\pi$ m $\approx 0{,}314$ m
\end{enumerate}
\end{document}
